% 全文统一的模型假设、题面条件与符号说明
\section{模型假设}

题面已经明确规定的区域尺度、频道数量、动作耗时与设备阈值均作为已知条件使用，
不再列作假设；本文只对题面没有完全刻画、且会影响后续推导的部分作如下统一假设。

\begin{itemize}[itemindent=2em]
  \item \textbf{假设 $1$（静态平面几何）：}一次任务内各干扰源位置保持不变，检测点的平面坐标已知且执行无误差，测向机与干扰源的高度差按题面规定忽略（依据：题面条款与定位必要理想化；用于问题一至问题四）。
  \item \textbf{假设 $2$（目标属性稳定）：}在题面规定的每频道至多一个干扰源的前提下，同一频道的全部有效观测对应同一干扰源，且该源的有效接收半径与定向方向在一局内保持不变（依据：题面条款与任务时段稳定性口径；用于问题一至问题四）。
  \item \textbf{假设 $3$（确定性误差解释）：}不同检测点的示向误差仅满足区间界而不预设独立性或概率分布，同一检测点的固定误差不通过重复检测取平均，所有保证性结论均按误差盒的最坏情形建立（依据：题面条款与最坏情形口径；用于问题一至问题四）。
  \item \textbf{假设 $4$（理想覆盖传播）：}除题面给定的距离衰减、定向覆盖半平面和示向误差外，不再引入额外遮挡、多径或源间干扰，覆盖区内场强只随距离变化（依据：题面条款与模型闭合需要；用于问题一至问题四）。
  \item \textbf{假设 $5$（理想运动控制）：}机器狗视为可即时改变行进方向的质点，指令坐标能够准确到达，除直线移动和题面列明的串行动作耗时外不计转向、加减速及通信计算时间（依据：题面直线运动规则与最短时间建模口径；用于问题二至问题四）。
  \item \textbf{假设 $6$（区域性质与空集处理）：}定位区域按连续闭集处理；半平面交集为空时判为观测不相容或候选构型不可行，不把空集解释为零面积；区域退化为单点或线段时按退化集合定义直径（依据：区域性质与数值稳健性口径；用于问题一至问题四）。
\end{itemize}

\subsection{题面条件与统一口径}

为避免把已知常量误写成假设，四问共同使用的题面条件汇总如下。

\begin{center}
\begin{tabularx}{\textwidth}{L C X}
  \toprule
  条件 & 数值或取值 & 统一解释 \\
  \midrule
  目标圆域半径 & $R_0=1800\ \mathrm{m}$ & 约束干扰源位置，不限制机器狗检测点位于圆域内 \\
  干扰源与频道 & $10\le N_J\le16$，频道编号为 $1$--$20$ & 问题三、四中每频道至多一个源，各源频道互异 \\
  示向误差 & $|e_i|\le\delta=1^\circ$ & 题面保证界，无概率分布；示向度统一按逆时针方位角定义 \\
  有效接收半径 & $1000\le R_j^{\mathrm{recv}}\le1500\ \mathrm{m}$ & 半径不可查询；保证接收按下限，保证超距无信号按上限 \\
  定向覆盖 & 定向方向两侧各 $90^\circ$（含边界） & 几何上表示包含边界的闭半平面，覆盖区内场强与角度无关 \\
  近距离检测 & 距离不超过 $5\ \mathrm{m}$ & 位于有效覆盖区时返回“距离过近”，可直接进入清除 \\
  光学定位与清除 & 距离不超过 $20\ \mathrm{m}$ & 对指定频道的未清除源，成功只取决于距离，与干扰源朝向无关 \\
  运动与检测 & 速度 $5\ \mathrm{m/s}$，检测 $5\ \mathrm{s}$，换频 $1\ \mathrm{s}$ & 换频只由相邻两次合法检测的频道变化触发 \\
  清除耗时 & 成功 $5\ \mathrm{s}$，未发现 $3\ \mathrm{s}$ & \texttt{/clear} 不切换测向机当前频道 \\
  \bottomrule
\end{tabularx}
\end{center}

数学模型使用题面误差界 $\delta=1^\circ$；在线程序另取
$\delta_{\mathrm{num}}=1.01^\circ$，其中增加的 $0.01^\circ$ 只用于吸收示向度两位小数表示和浮点运算的量化余量，
不改变问题一、二的物理误差参数，也不把题面保证值改写为 $1.01^\circ$。
问题一、二只对接口返回正常示向度的观测建模；该响应同时给出
$5<\lVert G-S_i\rVert_2\le R_j^{\mathrm{recv}}\le1500\ \mathrm{m}$，首次正检测还给出
$R_j^{\mathrm{recv}}\ge\max\{1000,\lVert G-S_1\rVert_2\}$，后者用于构造问题二的条件保证接收区。
问题二按题面指定的全向干扰源建立站址模型，不把问题四的定向盲区并入第二检测点可行域。
圆域、闭测向锥及其交集均按闭凸集处理；交集为空或测向 Jacobian 退化时把候选构型判为不可行，
交集退化为单点或线段时则分别按直径为零或线段长度计算。
问题一用正 $720$ 边形内、外逼近目标圆域并同时报告直径下界和上界；在线定位则用外接正 $360$ 边形包围目标圆域，
并用 $128$ 个圆的支撑半平面外包有效接收圆，故用于清除的可行集始终是保守外包围。

数值实现把物理边界向安全侧内缩：保证接收距离采用 $999.99\ \mathrm{m}$，
问题四连续域发现证书采用 $994.99\ \mathrm{m}$，保证超距无信号的程序判据采用 $1500.01\ \mathrm{m}$。

认证清除半径采用 $19.999\ \mathrm{m}$，移动终点的剩余半径采用 $19.998\ \mathrm{m}$，覆盖单元半径采用 $19.5\ \mathrm{m}$；
半平面裁剪的基本距离保护量为 $10^{-6}\ \mathrm{m}$，相同站位判定阈值为 $10^{-7}\ \mathrm{m}$，
最坏测向包围另加 $10^{-3}\ \mathrm{m}$ 外扩量；通用几何比较采用 $10^{-10}$ 的相对容差，矩阵秩判断采用 $10^{-12}$ 的相对阈值，交会角退化判断采用 $10^{-12}$ 的无量纲正弦阈值，这些量均为工程余量而非新的物理参数。

对尚未清除的频道，问题三中的 \texttt{no\_signal} 可能表示“该频道无源”或“超出该源的有效接收半径”；
在条件“该频道确有全向源”下才只剩距离原因。
问题四还增加“检测点位于定向盲区”这一原因，因而只有扫描点的局部凸包证书覆盖整个目标圆域后，
才能把从未获得正检测的频道记为不存在。
若已经发现 $16$ 个异频干扰源，则由源数上界可直接认证其余频道无源；除此以外，
任务完成的程序判据是全部 $20$ 个频道均处于“已清除”或“确认不存在”状态。

程序严格串行调用 \texttt{/enter}、\texttt{/measure}、\texttt{/clear} 与 \texttt{/exit}；
每个新动作使用新的 \texttt{request\_id}，网络故障只以完全相同的请求体和同一标识重试，
因此幂等重试不会增加有效动作或虚拟时间。
题面没有规定 \texttt{/clear} 返回“未发现”的次数上限；问题四在线程序设置每局最多 $5$ 次的工程预算，
预算用尽后转入有误差上界的测向包，或仅在“既有失败圆与当前圆覆盖全部剩余可行域”的末单元证书成立时清除，
所以该预算限制试探代价而不削弱完备性。
程序不读取案例真值，只以四个公开接口的合法响应更新状态。

\section{符号说明}

下表只保留四问正文及统一时间模型所需的符号；“首次出现”指其首次参与实质模型推导的章节。

\begin{center}
\begin{tabularx}{\textwidth}{C X C C}
  \toprule
  符号 & 含义 & 单位或取值 & 首次出现 \\
  \midrule
  $S_i=(x_i,y_i)^{\mathsf T}$ & 第 $i$ 个检测点的位置向量 & $\mathrm{m}$ & 问题一 \\
  $G_j=(x_j,y_j)^{\mathsf T}$ & 第 $j$ 个干扰源的位置向量；单源推导时简写为 $G$ & $\mathrm{m}$ & 问题一 \\
  $\hat\theta_i$ & 第 $i$ 个检测点处测得的示向度 & $[0^\circ,360^\circ)$ & 问题一 \\
  $\theta_i$ & 从检测点 $S_i$ 指向干扰源的真实方位角 & $[0^\circ,360^\circ)$ & 问题一 \\
  $e_i$ & 第 $i$ 次检测的示向误差，$\hat\theta_i=\theta_i+e_i$ & $({}^\circ)$ & 问题一 \\
  $\delta$ & 题面给定的示向误差上界 & $1^\circ$ & 问题一 \\
  $\delta_{\mathrm{num}}$ & 在线定位使用的示向误差半角 & $1.01^\circ$ & 问题三 \\
  $u(\phi)$ & 方位角 $\phi$ 对应的二维单位方向向量 & 无量纲 & 问题一 \\
  $W_i$ & 第 $i$ 次示向度对应的闭凸测向锥 & 集合 & 问题一 \\
  $K,\ K_q^-,\ K_q^+$ & 精确定位区域及其内接、外接正 $q$ 边形逼近 & 集合 & 问题一 \\
  $q$ & 目标圆域离散逼近的正多边形边数 & $720$ & 问题一 \\
  $D(K)$ & 定位区域内任意两点距离的最大值 & $\mathrm{m}$ & 问题一 \\
  $r_{\mathrm{MEC}}(K)$ & 定位区域最小覆盖圆的半径 & $\mathrm{m}$ & 问题一 \\
  \bottomrule
\end{tabularx}
\end{center}

\begin{center}
\begin{tabularx}{\textwidth}{C X C C}
  \toprule
  符号 & 含义 & 单位或取值 & 首次出现 \\
  \midrule
  $R_0$ & 目标圆域半径 & $1800\ \mathrm{m}$ & 问题一 \\
  $R_j^{\mathrm{recv}}$ & 第 $j$ 个干扰源的有效接收半径 & $[1000,1500]\ \mathrm{m}$ & 问题二 \\
  $U_1$ & 第一次正常示向度给出的完整可能源集合 & 集合 & 问题二 \\
  $X,\ S_2$ & 第二检测点候选位置与最终选定位置 & $\mathrm{m}$ & 问题二 \\
  $\mathcal C_{1000},\ \mathcal C_{\mathrm{det}}$ & 保守保证接收区与利用首次正检测信息的条件保证区 & 集合 & 问题二 \\
  $J$ & 测向模型对干扰源位置的 Jacobian 矩阵 & $\mathrm{rad/m}$ & 问题二 \\
  $K_{\mathrm{LS}}$ & 测角残差到位置增量的最小二乘增益矩阵 & $\mathrm{m/rad}$ & 问题二 \\
  $\Sigma_G$ & 一阶位置估计协方差矩阵 & $\mathrm{m^2}$ & 问题二 \\
  $\gamma$ & 两条交会视线的锐夹角 & $[0^\circ,90^\circ]$ & 问题二 \\
  $\mathrm{GDOP}$ & 测向几何精度衰减因子 & $\mathrm{m/rad}$ & 问题二 \\
  $f_1(X),\ f_2(X)$ & 确定性最坏位置误差与第二次检测耗时 & $\mathrm{m},\ \mathrm{s}$ & 问题二 \\
  $\mathcal P$ & 精度与耗时的 Pareto 非支配候选点集 & 集合 & 问题二 \\
  \bottomrule
\end{tabularx}
\end{center}

\begin{center}
\begin{tabularx}{\textwidth}{C X C C}
  \toprule
  符号 & 含义 & 单位或取值 & 首次出现 \\
  \midrule
  $P_k$ & 第 $k$ 个扫描点或路径节点的位置 & $\mathrm{m}$ & 问题三 \\
  $N_J$ & 一局中的干扰源总数 & $10,\ldots,16$ & 问题三 \\
  $N_m,N_s$ & 检测次数与换频次数 & 次 & 问题三 \\
  $N_c^{\mathrm{succ}},N_c^{\mathrm{miss}}$ & 成功清除与未发现清除的次数 & 次 & 问题三 \\
  $T$ & 定位清除任务的虚拟总时间 & $\mathrm{s}$ & 问题三 \\
  $V_G$ & 距可能源位置 $G$ 不超过 $1000\ \mathrm{m}$ 的扫描点相对向量集 & 集合 & 问题四 \\
  $d$ & 定向干扰源的单位发射方向向量 & 无量纲 & 问题四 \\
  $P_4,\ P_4'$ & 问题四的基准扫描点集与连续域认证提速点集 & 集合 & 问题四 \\
  \bottomrule
\end{tabularx}
\end{center}

\begin{table}[H]
\centering
\begin{minipage}[t]{0.9\textwidth}
\centering
\caption{术语与状态表}
\small
\begin{tabular}{@{}p{3.4cm}p{9.6cm}@{}}
\toprule
术语 & 含义（单位） \\
\midrule
MEC & 最小覆盖圆 (Minimum Enclosing Circle) \\
GDOP & 几何精度衰减因子 $\sqrt{\operatorname{tr}[(H^{\mathsf T}H)^{-1}]}$ (m/rad) \\
RMS & 一阶位置均方根误差 (m) \\
Pareto 前沿 & 两目标下不被其他可行解支配的候选点集 (集合) \\
\texttt{unknown} & 尚未发现该频道源 (状态枚举) \\
\texttt{positively\_detected} & 该频道至少得到一次正常示向度或“距离过近” (状态枚举) \\
\texttt{localizing} & 正在追加异地观测并收缩定位区域 (状态枚举) \\
\texttt{clear\_ready} & 外近似 MEC 半径不超过 20 m，或返回“距离过近” (状态枚举) \\
\texttt{cleared} & 该频道干扰源已成功清除 (状态枚举) \\
\texttt{absent\_certified} & 保证扫描全部完成且该频道从未被正检测 (状态枚举) \\
\texttt{geometry\_fault} & 保证分支失信号、空交集或认证清除失败 (阻断状态) \\
\texttt{clear\_fallback} & 对当前外近似执行有限 20 m 覆盖清除 (状态枚举) \\
\texttt{request\_id} & 单次测试会话内的 HTTP 幂等键 (字符串) \\
虚拟时间 & 模拟器按移动与动作规则累计的任务耗时 (s) \\
现实截止时刻 & \texttt{/enter} 返回的剩余程序运行时间约束 (s) \\
\bottomrule
\end{tabular}
\end{minipage}
\end{table}

\section{问题分析}

\textbf{问题一：}此问要求根据检测点坐标和示向度确定定位区域并求其直径，进而判断直径圆能否覆盖该区域。由于测向误差不超过 \(1^\circ\)，不能将示向度视为精确射线，而应把真实方位角视为落在以测得方向为中心、张角 \(2^\circ\) 的扇形内。因此，每个检测点给出一个凸锥约束，干扰源必须位于所有凸锥的交集内。同时，干扰源还受限于半径 \(1800\,\mathrm{m}\) 的目标圆域，故定位区域为多个凸锥与圆域的交集。若交集为空，说明观测不相容；若退化为点或线段，则按退化情形处理。求解时，可将圆域用多边形逼近，把问题转化为半平面求交，得到凸多边形定位区域，再利用凸性求其直径。最后，针对直径圆能否覆盖区域，需分析区域形状，并考虑反例与最小覆盖圆的关系。整体思路是：误差建模 \(\to\) 凸锥交会 \(\to\) 圆域约束 \(\to\) 多边形逼近 \(\to\) 直径计算 \(\to\) 覆盖性讨论。

\textbf{问题二：}此问基于第一问覆盖问题的求解，考虑第一次检测点 \(S_1\) 及示向度 \(\hat\theta_1\) 已知，而第二次检测点 \(S_2\) 待设计的情况。第一次观测给出带误差的示向锥，结合目标距离大于 \(5\mathrm{m}\)、接收半径不超过 \(1500\mathrm{m}\) 及目标圆域，得可能源集合 \(U_1\)。第二检测点须先保证能再次收到信号：最保守取到 \(U_1\) 任意点距离不超过 \(1000\mathrm{m}\) 的核心保证区 \(\mathcal C_{1000}\)；利用第一次成功接收提供的距离信息，可放宽为条件保证区 \(\mathcal C_{\det}\)。精度目标是两次交会后定位区域面积最小，但第二次示向度未知且面积非凸，故采用一阶误差传播构造面积代理 \(A_{\mathrm{cert}}=\pi F_1^2\)，选点后再用问题一模型计算实际区域与最小覆盖圆。按“\(U_1\)—候选区—面积代理—Pareto 选点—数值验证”展开，最终给出第二检测点可行分布与推荐策略。接收可靠是首要约束，不能只追求交会角接近直角。

\textbf{问题三：}此问在前面两问研究的基础上，要求机器狗从原点出发、初始频道为 1，在 10–16 个全向干扰源数量未知、频道未知、位置未知、有效接收半径未知的条件下，完成全部干扰源的定位与清除，并使总时间尽可能短。与问题一、二相比，本问的核心困难有三点：1.发现性：没有任何先验位置信息，必须设计一套有限步内必能发现全部干扰源的扫描方案，并且要能对"某频道不存在"给出结论，否则无法判断任务是否真正完成。
2.定位的收敛性：每次测向只给出 $1^\circ$ 的角度信息，单次交会无法直接达到光学清除所需的精度，需要设计保证能持续获得新观测的取点规则。
3.计时口径：总时间包含移动、换频、检测、光学精确定位与清除，各项耗时规则不同（移动按距离除以 5 m/s，换频 1 s，检测 5 s，清除成功 5 s、未发现 3 s），策略必须与实际计时规则一一对应。

\textbf{问题四：}此问在问题三的基础上，新增要求：目标区域内除全向干扰源外还存在定向干扰源，其个数与定向方向均未知。这一差别改变了问题的一个基本逻辑，即问题三中"某个频道扫不到信号"只有两种解释：该频道没有干扰源，或距离超过有效接收半径，两者都能由扫描覆盖证书排除，所以"扫描点全部测完仍无信号"可以作为"该频道不存在"的证明。而问题四中多出第三种解释：检测点位于该定向干扰源的盲区。定向干扰源的有效覆盖角为定向方向两侧各 $90^\circ$，覆盖角之外没有信号辐射。为解决新增情况带来的判断问题，先设计对任意未知朝向都保证发现的扫描骨架。固定可能源位置 \(G\)，记其接收半径内扫描点为 \(V_G\)。若全部扫描点均不可见，则存在方向使所有 \(P_j-G\) 落在同一开半平面内，等价于 \(0\notin\operatorname{conv}V_G\)。故只要保证
\[
G\in\operatorname{conv}\{P_j:\|P_j-G\|\le1000\},\quad \forall G\in B_{1800},
\]
就能对任意朝向至少有一个扫描点可见。据此可取边长 \(990\,\mathrm{m}\) 的等边三角格，扫描其与目标圆域相交的基本三角形顶点，共 \(31\) 点；或用 \(22\) 点认证集配合连续域凸包证书提速。
同时在定位与清除上，首个示向度仍给出位置扇形，但首点可能位于盲区边界，后续 \texttt{no\_signal} 不再视为故障，而是触发完备路径：在扇形上铺 \(20\,\mathrm{m}\) 三角格逐点清除，有限步内必清除成功；一旦新增观测使最小覆盖圆半径降到 \(20\,\mathrm{m}\) 内，则回到单点清除。状态转移与问题三基本相同，但“确认不存在”必须满足上述凸包条件。

\subsection{总体框架}
